\documentclass[a4paper,12pt]{article}

\usepackage[T2A]{fontenc}
\usepackage[utf8]{inputenc}
\usepackage[russian]{babel}
\usepackage{amsmath,amssymb,mathtools}
\usepackage{PTSans}
\renewcommand{\familydefault}{\sfdefault}
\usepackage{icomma}
\usepackage{geometry}
\usepackage{microtype}
\usepackage{tikz}
\usetikzlibrary{calc,arrows.meta,positioning,shapes.geometric}
\usepackage{tabularx,booktabs,array}
\usepackage{enumitem}
\usepackage{fancyhdr}
\usepackage{setspace}
\usepackage{xcolor}
\usepackage{hyperref}
\usepackage{parskip}

\geometry{left=20mm,right=20mm,top=18mm,bottom=18mm}
\onehalfspacing
\setlength{\parindent}{0pt}
\setlength{\parskip}{0.55em}
\setlist[itemize]{leftmargin=1.35em,itemsep=0.35em,topsep=0.3em}
\setlist[enumerate]{leftmargin=1.65em,itemsep=0.45em,topsep=0.3em}

\definecolor{accent}{HTML}{008080}
\definecolor{darkaccent}{HTML}{004D4D}
\definecolor{textgray}{HTML}{333333}
\definecolor{softgray}{HTML}{F2F5F5}

\hypersetup{hidelinks}
\color{textgray}

\pagestyle{fancy}
\fancyhf{}
\fancyfoot[C]{\thepage}
\renewcommand{\headrulewidth}{0pt}
\renewcommand{\footrulewidth}{0pt}

\newcommand{\sectiontitle}[1]{%
  \vspace{0.4em}%
  {\Large\bfseries\color{darkaccent}#1}\par
  \vspace{0.35em}%
  {\color{accent}\rule{\linewidth}{0.7pt}}\par
  \vspace{0.5em}%
}

\tikzset{
  startstop/.style={ellipse, draw=darkaccent, very thick, align=center, fill=softgray,
    minimum height=11mm, text width=8.2cm, inner sep=4pt},
  process/.style={rounded corners=3pt, draw=accent, thick, align=center, fill=white,
    minimum height=12mm, text width=8.7cm, inner sep=5pt},
  decision/.style={diamond, aspect=2.1, draw=darkaccent, thick, align=center, fill=softgray,
    inner sep=2pt, text width=5.0cm},
  arrow/.style={-{Stealth[length=3mm]}, thick, draw=darkaccent}
}

\begin{document}

% -------------------- Титульный лист --------------------
\thispagestyle{empty}
\begin{center}
  \vspace*{2.6cm}

  {\fontsize{22}{27}\selectfont\bfseries\color{darkaccent} Чек-лист}\par
  \vspace{1.1cm}

  {\fontsize{25}{31}\selectfont\bfseries Линейные уравнения: от текста к уравнению}\par
  \vspace{0.55cm}

  {\large Решение уравнений, раскрытие скобок, подобные слагаемые и перевод\\
  словесных условий}\par

  \vfill

  {\Large\bfseries Кирилл Зиновьев}\par
  \vspace{0.7cm}
  {\large 30.09.26}\par

  \vfill

  {\normalsize Лёвин Артём Александрович эксклюзивно для Кирилла Зиновьева}\par
  \vspace{0.8cm}

  \begin{tikzpicture}[x=1cm,y=1cm]
    \draw[accent, line width=1.2pt]
      (-6,0) .. controls (-3.5,0.75) and (-1.3,-0.65) .. (0,0)
             .. controls (1.6,0.8) and (4.2,-0.6) .. (6,0.15);
  \end{tikzpicture}
\end{center}
\clearpage

% -------------------- Основной материал --------------------
\sectiontitle{1. Что уже должно получаться уверенно}

\begin{itemize}
  \item раскрывать скобки с учётом знака перед ними;
  \item приводить подобные слагаемые;
  \item переносить слагаемые из одной части уравнения в другую с изменением знака;
  \item делить обе части уравнения на ненулевой коэффициент при неизвестной;
  \item проверять найденное значение подстановкой.
\end{itemize}

\textbf{Пример.}
\[
-2(x-1)+3(4-x)=10.
\]
Раскрываем скобки и приводим подобные:
\[
-2x+2+12-3x=10,
\]
\[
-5x+14=10,
\qquad
-5x=-4,
\qquad
x=\frac45.
\]

Ещё один типичный вид:
\[
2x-3=-x+5
\quad\Longrightarrow\quad
3x=8
\quad\Longrightarrow\quad
x=\frac83.
\]

\sectiontitle{2. Главный новый навык: перевести текст на язык уравнений}

Перед записью уравнения полезно мысленно ответить на три вопроса:

\begin{enumerate}
  \item Что сравнивают?
  \item С чем сравнивают?
  \item Как именно сравнивают: на сколько или во сколько раз?
\end{enumerate}

\renewcommand{\arraystretch}{1.45}
\begin{center}
\begin{tabularx}{\linewidth}{>{\raggedright\arraybackslash}X >{\centering\arraybackslash}p{4.4cm} >{\raggedright\arraybackslash}X}
\toprule
\textbf{Фраза} & \textbf{Уравнение} & \textbf{Как понимать} \\
\midrule
$A$ больше $B$ на $k$ & $A=B+k$ & к меньшему прибавили $k$ \\
$A$ меньше $B$ на $k$ & $A+k=B$ или $A=B-k$ & разность равна $k$ \\
$A$ больше $B$ в $k$ раз & $A=kB$ & большее равно $k$ меньшим \\
$A$ меньше $B$ в $k$ раз & $A=\dfrac{B}{k}$ или $kA=B$ & меньшее равно одной $k$-й от большего \\
\bottomrule
\end{tabularx}
\end{center}

\textbf{Важно.} Слова «на» и «в ... раз» задают разные действия:
\[
\text{«на }k\text{»} \Rightarrow \text{сложение/вычитание},
\qquad
\text{«в }k\text{ раз»} \Rightarrow \text{умножение/деление}.
\]

\clearpage
\sectiontitle{3. Примеры перевода формулировок}

\begin{enumerate}
  \item «При каком значении $b$ выражение $8b-27$ равно $5$?»
  \[
  8b-27=5.
  \]
  Отсюда $b=4$.

  \item «Значения выражений $2m-13$ и $m-3$ равны».
  \[
  2m-13=m-3.
  \]
  Отсюда $m=10$.

  \item «$3-5c$ на $1$ меньше, чем $1-c$».
  \[
  3-5c=(1-c)-1.
  \]
  Отсюда $c=\dfrac34$.

  \item «$2x+1$ на $20$ больше, чем $8x+5$».
  \[
  2x+1=(8x+5)+20.
  \]
  Отсюда $x=-4$.

  \item «$x$ в $3$ раза меньше, чем $45-10x$».
  \[
  x=\frac{45-10x}{3}
  \quad\Longleftrightarrow\quad
  3x=45-10x.
  \]
  Отсюда $x=\dfrac{45}{13}$.

  \item «$9-y$ в $2$ раза больше, чем $y$».
  \[
  9-y=2y.
  \]
  Отсюда $y=3$.
\end{enumerate}

\sectiontitle{4. Частые ошибки}

\begin{itemize}
  \item Добавочное число ставят не к той части. Если $A$ больше $B$ на $10$, то $A=B+10$.
  \item Путают «на» и «в раз». «На 5» связано с $+5$ или $-5$; «в 5 раз» — с умножением или делением на $5$.
  \item Сразу переносят числа, не поняв смысл текста. Сначала составьте правильное уравнение, затем решайте его.
  \item Теряют знак при раскрытии скобок. Например, $-2(x-1)=-2x+2$.
  \item Не выполняют проверку. Подстановка найденного значения должна делать исходное равенство верным.
\end{itemize}

\clearpage
% -------------------- Блок-схема --------------------
\thispagestyle{empty}
\begin{center}
{\Large\bfseries\color{darkaccent} Алгоритм: как перевести текст в уравнение}\par
\vspace{1cm}

\begin{tikzpicture}[node distance=1.05cm and 1cm]
  \node[startstop] (start) {Прочитать условие полностью};
  \node[process, below=of start] (objects) {Определить два сравниваемых выражения:\\что сравнивают и с чем};
  \node[decision, below=1.2cm of objects] (kind) {В условии сказано\\«на» или «в ... раз»?};

  \node[process, below left=1.7cm and 0.65cm of kind, text width=5.25cm] (add) {Если «на $k$»: использовать сложение или вычитание.\\Например: $A$ больше $B$ на $k$ $\Rightarrow A=B+k$};
  \node[process, below right=1.7cm and 0.65cm of kind, text width=5.25cm] (mult) {Если «в $k$ раз»: использовать умножение или деление.\\Например: $A$ меньше $B$ в $k$ раз $\Rightarrow kA=B$};
  \path (add.south) -- (mult.south) coordinate[midway] (branchesmid);
  \node[process, below=1.55cm of branchesmid] (equation) {Записать равенство и проверить его смысл словами};
  \node[process, below=of equation] (solve) {Если требуется: решить полученное линейное уравнение};
  \node[startstop, below=of solve] (check) {Подставить ответ и проверить исходное условие};

  \draw[arrow] (start) -- (objects);
  \draw[arrow] (objects) -- (kind);
  \draw[arrow] (kind) -- node[above left, font=\small]{«на»} (add);
  \draw[arrow] (kind) -- node[above right, font=\small]{«в ... раз»} (mult);
  \draw[arrow] (add.south) -- ++(0,-0.45) -| ([xshift=-2.6cm]equation.north);
  \draw[arrow] (mult.south) -- ++(0,-0.45) -| ([xshift=2.6cm]equation.north);
  \draw[arrow] (equation) -- (solve);
  \draw[arrow] (solve) -- (check);
\end{tikzpicture}
\end{center}
\clearpage

% -------------------- Тренировочный блок --------------------
\sectiontitle{Тренировочный блок — Путь А}

\textbf{10 заданий.} В задачах 1–3 решите уравнение. В задачах 4–10 сначала составьте уравнение по тексту, затем решите его.

\begin{enumerate}
  \item $4x-7=13$.
  \item $-3(x-2)+2(5-x)=11$.
  \item $5x-4=-2x+17$.
  \item При каком значении $a$ выражение $6a-11$ равно $19$?
  \item При каком значении $m$ значения выражений $3m-8$ и $m+6$ равны?
  \item Значение выражения $4x+3$ на $11$ больше значения выражения $2x-5$. Найдите $x$.
  \item Значение выражения $7y-9$ на $6$ меньше значения выражения $3y+13$. Найдите $y$.
  \item Значение выражения $5p-1$ в $3$ раза больше значения выражения $p+7$. Найдите $p$.
  \item Значение выражения $q+4$ в $4$ раза меньше значения выражения $5q-8$. Найдите $q$.
  \item Значение выражения $2z-3$ на $15$ больше значения выражения $8-3z$. Найдите $z$.
\end{enumerate}

\vfill
\textbf{Ориентир:} сначала установите отношение между двумя выражениями, только затем переходите к преобразованиям.
\clearpage

% -------------------- Ответы --------------------
\sectiontitle{Ответы}

\renewcommand{\arraystretch}{1.45}
\begin{center}
\begin{tabularx}{\linewidth}{>{\centering\arraybackslash}p{0.9cm} >{\raggedright\arraybackslash}X >{\centering\arraybackslash}p{3.1cm}}
\toprule
\textbf{№} & \textbf{Ключевое уравнение} & \textbf{Ответ} \\
\midrule
1 & $4x-7=13$ & $x=5$ \\
2 & $-3(x-2)+2(5-x)=11$ & $x=1$ \\
3 & $5x-4=-2x+17$ & $x=3$ \\
4 & $6a-11=19$ & $a=5$ \\
5 & $3m-8=m+6$ & $m=7$ \\
6 & $4x+3=(2x-5)+11$ & $x=\dfrac32$ \\
7 & $(7y-9)+6=3y+13$ & $y=4$ \\
8 & $5p-1=3(p+7)$ & $p=11$ \\
9 & $4(q+4)=5q-8$ & $q=24$ \\
10 & $2z-3=(8-3z)+15$ & $z=\dfrac{26}{5}$ \\
\bottomrule
\end{tabularx}
\end{center}

\vfill
\begin{center}
{\small Лёвин Артём Александрович эксклюзивно для Кирилла Зиновьева}
\end{center}

\end{document}
